% Laget automatisk av elevark/lag_elevark.py fra sim/romreise.html. Endringer her blir overskrevet.
\documentclass[11pt]{article}
\usepackage{elevark}
\begin{document}
\elevarkhode{Fysikk 2 · Spesiell relativitetsteori}{Romreise}{https://lektorbo.no/sim/romreise.html}{Simuleringen}

Alfa Centauri ligger 4,37~lysår unna. Ingenting kan gå fortere enn lyset, så sett fra Jorda tar reisen minst 4,37~år. Men for mannskapet kan reisen ta mye kortere tid. Sett fra romskipet er avstanden til stjerna kortere, fordi strekningen beveger seg forbi skipet med en fart nær lysfarten.

\section*{Oppgaver}

Start med Alfa Centauri og $0{,}90c$.

\begin{oppgaveliste}
\oppgave \fase{test}{Regn}Regn ut reisetiden sett fra Jorda. Regn så ut $γ$, og avstanden til Alfa Centauri i skipets system. Hvor lang tid tar reisen for mannskapet?
\svar
\oppgave \fase{forklar}{Forklar}På Jorda sier vi at reisen tar kortere tid for mannskapet fordi klokkene på skipet går saktere. Hvordan forklarer mannskapet selv at de kommer fram så fort? Bytt til skipets system.
\svar
\oppgave \fase{gjett}{Gjett}Hva skjer med reisetiden på skipet når farten nærmer seg lysfarten? Kan den bli null? Se på grafen.
\svar
\oppgave \fase{test}{Regn}Du vil til Vega (25~lysår) og ha brukt 10~år på skipsklokka når du kommer fram. Hvor fort må du reise? Hvor gammel er du ved ankomst, og hvor mange år har gått på Jorda?
\svar
\oppgave \fase{forklar}{Forklar}Med $0{,}99c$ til Alfa Centauri tar reisen 0,62~år på skipet. Mannskapet har reist 4,37~lysår på under ett år. Betyr det at de har reist fortere enn lyset?
\svar
\oppgave \fase{test}{Regn}Du reiser til Alfa Centauri og tilbake med $0{,}90c$. Tvillingen din blir igjen på Jorda. Hvor gamle er dere når du kommer hjem?
\svar
\oppgave \fase{utfordring}{Utfordring}Hvor fort må du reise for å nå sentrum av Melkeveien (26 000~lysår) på 40~år skipstid? Klarer du det med den høyeste farten i simuleringen?
\svar
\end{oppgaveliste}

\end{document}
